%----------
%	MÉTODOS DE DETECCIÓN DE SESGOS DE GÉNERO
%----------	

La detección y análisis de sesgos de género en el discurso mediático han experimentado una profunda transformación metodológica en la última década, impulsada por la necesidad de procesar volúmenes de información cada vez mayores. Conviven en ella dos grandes enfoques: el tradicional, basado en la codificación humana, y el automatizado, apoyado en técnicas de Inteligencia Artificial.


\subsection{Enfoques Tradicionales y Análisis Manual}

El análisis de la representación mediática y la publicación de los informes que se derivan de él descansan, todavía hoy, en un procedimiento manual. Los equipos de investigación y las expertas en género elaboran en primer lugar una metodología exhaustiva y detallada, materializada a menudo en libros de códigos extensos, como el que describe la tesis doctoral de Clara Sáinz de Baranda \cite{sainzdebaranda2013}. Después llega el trabajo de campo, que exige recorrer periódico a periódico y noticia a noticia, codificando cada pieza para el análisis cuantitativo, con ayuda de diversas herramientas y aplicaciones, y marcando el texto a mano para el cualitativo.


Este enfoque destaca por su alto nivel de precisión humana y comprensión del contexto sociocultural, siendo el estándar de referencia (\textit{golden standard} o \textit{ground truth}, en inglés) en investigaciones académicas. Como ya se ha mencionado, un ejemplo paradigmático de este rigor es la tesis doctoral de Clara Sáinz de Baranda \cite{sainzdebaranda2013}, que requirió el análisis manual de miles de piezas periodísticas para visibilizar la marginación de las mujeres en la prensa desde 1978 hasta 2010. Del mismo modo, proyectos recientes de gran envergadura como el ya mencionado InfoIA 2026 \cite{sainzdebaranda2026} siguen apoyándose en metodologías de revisión experta para garantizar la fiabilidad de los datos. Sin embargo, estas metodologías presentan limitaciones evidentes en términos de escalabilidad: conllevan una inmensa inversión de tiempo, una altísima carga de trabajo manual y la imposibilidad de auditar el ecosistema mediático en tiempo real.


\subsection{Transición hacia la Automatización e Inteligencia Artificial}
Para superar estas barreras de la escalabilidad manual, en los últimos años han surgido diversas herramientas computacionales y plataformas de monitorización que aplican técnicas de Procesamiento de Lenguaje Natural (PLN) e Inteligencia Artificial para la detección de sesgos y discursos discriminatorios. Estas iniciativas se pueden clasificar en tres grandes grupos según su complejidad técnica:

\begin{itemize}
    \item \textbf{Sistemas basados en diccionarios y glosarios:} Un primer acercamiento a la automatización en las redacciones es el uso de sistemas de validación léxica. Destaca el caso del diario argentino La Nación, que desarrolló la herramienta interna \textit{Genie} y su versión en abierto \textit{Cómo lo digo} \cite{lanacion2024}. Esta aplicación funciona cruzando el texto de la periodista con un glosario predefinido que detecta términos inapropiados o que, por su contexto, podrían perpetuar estereotipos. La herramienta marca estas expresiones y sugiere ``Palabras que se recomienda reemplazar'' o ``Palabras que se recomienda revisar'', actuando como un corrector ortotipográfico con perspectiva de género.

    \item \textbf{Monitorización del discurso:} A nivel macroscópico, existen rastreadores automatizados que auditan la producción periodística de forma masiva. El \textit{GenderGap Tracker}, desarrollado por Informed Opinions \cite{gendergap2018}, analiza diariamente los principales medios de comunicación de Canadá para cuantificar la proporción de voces femeninas frente a masculinas en las noticias. En el ámbito hispanohablante destaca \textit{Hatemedia} \cite{hatemedia2023}, un monitor español de odio en la red. Sin embargo, el hito institucional más relevante en este ámbito es la reciente puesta en marcha de \textit{HODIO} (Huella del Odio y la Polarización) \cite{hodio2026}. Impulsada por el Observatorio Español del Racismo y la Xenofobia (OBERAXE) del Ministerio de Inclusión, Seguridad Social y Migraciones, esta plataforma constituye un esfuerzo pionero desde la administración pública. \textit{HODIO} emplea algoritmos de Inteligencia Artificial para auditar masivamente las principales redes sociales (Instagram, TikTok, X, YouTube y Facebook), para clasificar, cuantificar y monitorizar la prevalencia y evolución de discursos discriminatorios y polarizantes.

    \item \textbf{Modelos avanzados de mitigación y postedición interactiva:} En el punto más avanzado de la investigación se encuentran las herramientas diseñadas no solo para detectar, sino para auditar a los propios algoritmos y ofrecer alternativas de forma interactiva. El proyecto \textit{RECAST} \cite{wright2021} es un sistema interactivo que permite visualizar las predicciones de los modelos de detección de toxicidad y proporciona sugerencias alternativas para el lenguaje marcado como tóxico, facilitando la interpretabilidad y la intervención del usuario sobre el texto.
\end{itemize}

Con todo, la mayoría de estas herramientas operan sobre reglas léxicas o clasificadores cerrados, con una capacidad limitada para razonar sobre el contexto o para justificar sus decisiones. La irrupción de los grandes modelos de lenguaje y de las arquitecturas de agentes abre una vía distinta, que se examina a continuación.